\section{LLM-to-KG Methods: Construction, Alignment, Enrichment, and Completion}

LLMs lower the cost of graph creation by replacing task-specific extraction models with schema-guided prompting and by generating descriptions or candidate relations. Their fluent output also makes malformed or unsupported structure easy to miss. The engineering question is not whether an LLM can emit triples. It is how candidate triples get normalized, qualified, validated, versioned, and promoted.

\subsection{Schema-constrained extraction and semantic normalization}

SPIRES is an early and influential pattern. A user supplies a LinkML schema, the system recursively prompts an LLM to fill nested fields, and it then grounds extracted values to existing ontologies through external normalization tools. It supports diverse biomedical schemas without task-specific training and is implemented in OntoGPT~\cite{caufield2024spires}. The separation between generative extraction and ontology grounding is the key idea, because it accepts that an LLM may recognize a relation linguistically while failing to assign a stable identifier.

Arsenyan et al.~\cite{arsenyan2024emr} tested twelve language-model architectures on extracting diseases, factors, treatments, and manifestations from EMR notes, scoring both structured output and hallucination against manual annotations. Bigger decoder-only models did not dominate. Extraction is still a constrained prediction problem, and smaller supervised models can be competitive, cheaper, and easier to validate.

RELATE handles what happens after free-form relation extraction. It embeds an ontology's predicates, retrieves candidate Biolink relations with SapBERT, and uses an LLM to rerank them while handling negation. On ChemProt it reported 52\% exact match but 94\% accuracy at rank ten, which shows candidate retrieval was much stronger than final relation selection~\cite{olasunkanmi2026relate}. That gap is diagnostically useful, because it separates ontology coverage and embedding retrieval from LLM disambiguation. Construction studies should report entity detection, identifier grounding, predicate candidate recall, final predicate accuracy, qualifier extraction, and negation as separate stages.

Schema constraints solve structural validity, not factual validity. A perfectly valid triple can reverse subject and object, erase negation, omit species, or turn correlation into causation. Promotion into an asserted biomedical KG should therefore require evidence spans, source identifiers, confidence calibration, and relation-specific validation rules.

\subsection{Human-in-the-loop construction at scale}

The Mental Disorders Knowledge Graph (MDKG) is the strongest recent example of an LLM embedded in a larger curation pipeline rather than used as an autonomous extractor. The workflow combines GPT-4-assisted pre-labeling, active-learning annotation, named-entity recognition, relation extraction, database parsing, and LLM-based triple refinement. Its relations retain conditional statements, cohort characteristics, provenance, and confidence, and the published graph contains more than ten million relations, including nearly one million associations reported as absent from existing resources~\cite{gao2025mdkg}.

MDKG demonstrates three practices that should become standard. Use the LLM to reduce annotation effort, not to replace independent validation. Attach contextual metadata to the relation rather than forcing every finding into an unconditional triple. Evaluate downstream utility separately from extraction quality. Its scale also carries a caution: ``novel'' means not found in the comparison resources, not experimentally validated. New graph edges are hypotheses or extracted claims until corroborated.

\subsection{LLM-assisted patient and task-specific graph construction}

GraphCare uses LLM output together with external biomedical KGs to assemble patient-specific concept graphs, then applies a bi-attention GNN for mortality, readmission, length-of-stay, and medication prediction. It reported large gains on MIMIC-III and MIMIC-IV, especially in low-data settings~\cite{jiang2024graphcare}. Here the LLM acts as an open-world knowledge source for graph construction, while the downstream predictor is not itself an LLM. The case belongs in the survey because it shows how generated graph content can shape a clinical model even without natural-language generation at inference.

Patient-specific graphs improve relevance and complicate validation. The graph depends jointly on coding quality, patient-history completeness, LLM extraction, external KG coverage, and graph composition rules. Random train/test splits can leak recurring codes or graph neighborhoods. Temporal patient splits, source-specific ablations, and audits of generated edges are needed before a gain can be attributed to personalization.

\subsection{Completion and semantic enrichment}

Completion methods use LLM semantics to predict missing relations or improve graph representations. DrKGC learns structural embeddings and logical rules with a lightweight model, retrieves a query-specific subgraph bottom-up, refines it with a GCN adapter, and integrates the structural signal into LLM fine-tuning. It reported improvements on general and biomedical completion datasets~\cite{xiao2025drkgc}. Its dynamic retrieval is more structure-aware than converting a fixed neighborhood to text, but predicted triples remain inductive outputs. Evaluation on randomly removed edges can overstate scientific discovery because neighboring structure may encode the answer, so temporal and source-held-out splits are stronger tests.

LLaDR prompts an LLM for treatment-oriented textual descriptions of biomedical concepts and uses those descriptions to fine-tune KG embeddings for drug repurposing~\cite{xiang2025lladr}. The approach can enrich sparse entities with mechanistic context that the KG does not represent explicitly. That ``common sense'' may also be memorized literature, unsupported inference, or outdated practice. Description provenance and counterfactual tests are essential, especially when a repurposing candidate is justified with text generated by the same model that helped rank it.

MedCo, discussed above, combines EHR-mined associations, type-constrained LLM relations, generated descriptions and rationales, and heterogeneous graph learning~\cite{kerdabadi2026medco}. Together, DrKGC, LLaDR, and MedCo mark a shift from using language only as node attributes toward using LLMs to synthesize relation semantics and training signals. They also make component attribution harder. At a minimum, studies should compare graph structure alone, verified source text alone, raw LLM descriptions, descriptions with source grounding, and the fused model.

\subsection{Construction quality as a release process}

The methods above motivate a four-state release model. A candidate is \emph{extracted} when the LLM emits a schema-valid triple linked to a source span. It becomes \emph{normalized} when entities and predicates map to controlled identifiers and qualifiers. It becomes \emph{verified} when automated constraints and an independent model or curator confirm the evidence mapping. It becomes \emph{asserted} only when an authorized curator approves it for a versioned production graph. Confidence scores do not replace these states; model probability, self-consistency, agreement among LLMs, source quality, and human approval measure different properties and should be stored separately. The graph must also support retraction, supersession, and rollback, or an automated updater can add claims without being able to correct them safely.
